\documentclass[12pt]{article}
\usepackage{ae}
\usepackage[T1]{fontenc}
\usepackage{amsmath}
\usepackage{amsfonts}
\usepackage{lmodern}
\usepackage[lmargin=1cm, rmargin=2cm,tmargin=2cm,bmargin=2cm]{geometry}
\usepackage{listings}
\usepackage{graphics,graphicx}
\usepackage{color}
\usepackage{xcolor}
\usepackage{float}
\usepackage{subcaption}
\usepackage{booktabs}
\author{\empty}

\title{Problem Set 1. Introduction to R \\ Quantitative Methods for Humanities}

\date{\empty}

\begin{document}
\maketitle

\begin{center}
    \textbf{\large 1. Understanding World Population Dynamics}    
\end{center}

Understanding population dynamics is crucial for many areas of social science. This exercise involves calculating basic demographic measures of births and deaths for the world population over two time periods: 1950–1955 and 2005–2010. The analysis will use data files \texttt{Kenya.csv}, \texttt{Sweden.csv}, and \texttt{World.csv}, which contain population data for Kenya, Sweden, and the world, respectively.

\subsection*{Data Description}

The variables included in the datasets are described in Table~\ref{tab:world-pop-data}.

\begin{table}[h]
\centering
\caption{Description of Variables in the Population Data}\label{tab:world-pop-data}
\begin{tabular}{ll}
\toprule
\textbf{Variable} & \textbf{Description} \\
\midrule
\texttt{country} & Abbreviated country name \\
\texttt{period} & Time period for data collection \\
\texttt{age} & Age group \\
\texttt{births} & Number of births (in thousands) \\
\texttt{deaths} & Number of deaths (in thousands) \\
\texttt{py.men} & Person-years for men (in thousands) \\
\texttt{py.women} & Person-years for women (in thousands) \\
\bottomrule
\end{tabular}
\end{table}

The data are collected for a five-year period, with person-years representing the time contribution of each individual during the period. For example, a person who lived through the entire five-year period contributes 5 person-years, while someone who lived only half the period contributes 2.5 person-years.

\begin{enumerate}
    \item \textbf{Crude Birth Rate (CBR):} The crude birth rate (CBR) is defined as:
    \[
    \text{CBR} = \frac{\text{Number of births}}{\text{Number of person-years lived}}.
    \]
    Compute the CBR for each period separately for Kenya, Sweden, and the world. 
    \begin{enumerate}
        \item Start by calculating the total person-years as the sum of person-years for men and women.
        \item Store the results as a vector of length 2 (one for each period) for each region with appropriate labels.
    \end{enumerate}
    Briefly describe any patterns observed in the resulting CBRs.

    \item \textbf{Total Fertility Rate (TFR):} The total fertility rate (TFR) adjusts for age composition in the female population. First, calculate the age-specific fertility rate (ASFR), defined as:
    \[
    \text{ASFR}[x, x+\delta) = \frac{\text{Number of births to women of age } [x, x+\delta)}{\text{Person-years lived by women of age } [x, x+\delta)}.
    \]
    Here, $x$ is the starting age, and $\delta$ is the width of the age range (e.g., 5 years). 
    \begin{enumerate}
        \item Compute the ASFR for Kenya, Sweden, and the world for each period.
        \item Using the ASFR, calculate the TFR as:
        \[
        \text{TFR} = \sum_{\text{Age Groups}} \text{ASFR}[x, x+\delta) \times \delta.
        \]
        Store the resulting TFRs for each region and period as vectors of length 2. Describe any patterns in reproduction across the regions and periods.
    \end{enumerate}

    \item \textbf{Crude Death Rate (CDR):} The crude death rate (CDR) is defined as:
    \[
    \text{CDR} = \frac{\text{Number of deaths}}{\text{Number of person-years lived}}.
    \]
    Compute the CDR for each period and region, and describe the patterns observed.

    \item \textbf{Age-Specific Death Rate (ASDR):} Calculate the age-specific death rate (ASDR) for each age group and region during the 2005–2010 period:
    \[
    \text{ASDR}[x, x+\delta) = \frac{\text{Number of deaths for people of age } [x, x+\delta)}{\text{Person-years of people of age } [x, x+\delta)}.
    \]
    Briefly describe the patterns observed in the ASDRs.

    \item \textbf{Counterfactual CDR:} To understand differences in the CDR between Kenya and Sweden, compute the counterfactual CDR for Kenya using Sweden’s age distribution:
    \[
    \text{CDR} = \sum_{\text{Age Groups}} \text{ASDR}[x, x+\delta) \times P[x, x+\delta),
    \]
    where $P[x, x+\delta)$ is the proportion of the population in the age range $[x, x+\delta)$, calculated as the ratio of person-years in that age range to the total person-years. Compare the counterfactual CDR with the observed CDR for Kenya, and provide a substantive interpretation.
\end{enumerate}

{\color{blue}
\begin{center}
    \textbf{\large Solution}    
\end{center}

\textbf{R code for analysis:}

\begin{verbatim}
## Load datasets
kenya <- read.csv("Kenya.csv")
sweden <- read.csv("Sweden.csv")
world <- read.csv("World.csv")

####
## Q1 - CBR
####

## Aggregate the person-year metric for both genders.
kenya$py <- kenya$py.men + kenya$py.women
sweden$py <- sweden$py.men + sweden$py.women
world$py <- world$py.men + world$py.women

## Explore the periods 
kenya$period; sweden$period; world$period
# OBSERVATION: there are 2 periods in all the datesets:
# 1. "1950-1955", ranging from observations 1 to 15
i1 <- 1:15 # saving the indexes of period 1
# 2. "2005-2010", ranging from observations 16 to 30
i2 <- 16:30 # saving the indexes of period 2
# More advanced
# i1 <- which(kenya$period == "1950-1955")
# i2 <- which(kenya$period == "2000-2005")
# Saving periods
periods <- unique(kenya$period)

## Create CBR in both periods
# Kenya’s
kc1 <- sum(kenya[i1, "births"]) / sum(kenya[i1, "py"])
kc2 <- sum(kenya[i2, "births"]) / sum(kenya[i2, "py"])
kCBR <- c(kc1, kc2)
names(kCBR) <- periods
# Sweden’s
sc1 <- sum(sweden[i1, "births"]) / sum(sweden[i1, "py"])
sc2 <- sum(sweden[i2, "births"]) / sum(sweden[i2, "py"])
sCBR <- c(sc1, sc2)
names(sCBR) <- periods
# World’s
wc1 <- sum(world[i1, "births"]) / sum(world[i1, "py"])
wc2 <- sum(world[i2, "births"]) / sum(world[i2, "py"])
wCBR <- c(wc1, wc2)
names(wCBR) <- periods
# Print CBRs
kCBR; sCBR; wCBR

## CONCLUSIONS: 
# 1. Kenya’s CBR declined from 0.052 to 0.039 between 1950–1955 and 2005–2010.
# 2. Sweden’s CBR also declined, but remained far more stable, 
#    going from 0.015 to 0.012 in between periods.
# 3. The world’s CBR also declined, from 0.037 to 0.020.

####
## Q2 - ASFR
####

## Explore the age groups
kenya$age; sweden$age; world$age
# Saving age groups
age_groups <- unique(kenya$age)

## Compute ASFR in both periods
# Kenya's
K.ASFR1 <- kenya$births[i1] / kenya$py.women[i1]
K.ASFR2 <- kenya$births[i2] / kenya$py.women[i2]
# Creating data
K.ASFR <- data.frame(row.names = age_groups, "1950-1955" = K.ASFR1, 
                     "2000-2005" = K.ASFR2, check.names = FALSE) 
# Sweden's
S.ASFR1 <- sweden[i1, "births"] / sweden[i1, "py.women"] # different way of subsetting
S.ASFR2 <- sweden[i2, "births"] / sweden[i2, "py.women"]
# Creating data
S.ASFR <- data.frame(row.names = age_groups, "1950-1955" = S.ASFR1, 
                     "2000-2005" = S.ASFR2, check.names = FALSE) 
# World's
W.ASFR1 <- world[i1, "births"] / world[i1, "py.women"]
W.ASFR2 <- world[i2, "births"] / world[i2, "py.women"]
# Creating data
W.ASFR <- data.frame(row.names = age_groups, "1950-1955" = W.ASFR1, 
                     "2000-2005" = W.ASFR2, check.names = FALSE) 

## Printing results
K.ASFR
S.ASFR
W.ASFR

## OBSERVATION:
# 1. From age groups 0-14 and 50-80+ (no fertile periods), there is no brith

## Compute TFR
K.TFR <- c(sum(K.ASFR1 * 5), sum(K.ASFR2 * 5))
S.TFR <- c(sum(S.ASFR1 * 5), sum(S.ASFR2 * 5))
W.TFR <- c(sum(W.ASFR1 * 5), sum(W.ASFR2 * 5))
# Naming
names(K.TFR) <- names(S.TFR) <- names(W.TFR) <- periods
# Print TFRs.
K.TFR; S.TFR; W.TFR

## OBSERVATIONS:
# 1. Kenya’s TFR declined significantly, from 7.6 to 4.9.
# 2. Sweden’s TFR decreased from 2.2 to 1.9.
# 3. The world’s TFR dropped by half, from 5.0 to 2.5.

####
## Q3 - CDR
####

## Compute the CDR
# Kenya's
K.CDR1 <- sum(kenya[i1, "deaths"]) / sum(kenya[i1, "py"])
K.CDR2 <- sum(kenya[i2, "deaths"]) / sum(kenya[i2, "py"])
K.CDR <- c(K.CDR1, K.CDR2)
# Sweden's
S.CDR1 <- sum(sweden[i1, "deaths"]) / sum(sweden[i1, "py"])
S.CDR2 <- sum(sweden[i2, "deaths"]) / sum(sweden[i2, "py"])
S.CDR <- c(S.CDR1, S.CDR2)
# World's
W.CDR1 <- sum(world[i1, "deaths"]) / sum(world[i1, "py"])
W.CDR2 <- sum(world[i2, "deaths"]) / sum(world[i2, "py"])
W.CDR <- c(W.CDR1, W.CDR2)
# Naming
names(K.CDR) <- names(S.CDR) <- names(W.CDR) <- periods 
## Print CDRs.
K.CDR; S.CDR; W.CDR

## OBSERVATIONS:
# 1. Kenya’s CDR decreased from 0.024 to 0.016.
# 2. Sweden’s CDR remained relatively constant at 0.01.
# 3. The world’s CDR fell from 0.019 to 0.008.
# 4. Puzzling finding: The CDR of Kenya in 2000-2005 is about the same as Sweden!
#    CDR ignores the age composition of a population. See next question!

####
## Q4 - ASDR
####

## Compute ASDR within the period 2000-2005
K.ASDR <- kenya$deaths[i2] / kenya$py[i2] # Kenya's
S.ASDR <- sweden$deaths[i2] / sweden$py[i2] # Sweden's
W.ASDR <- world$deaths[i2] / world$py[i2] # World's

## Creating data
ASDR <- data.frame(row.names = age_groups, Kenya = K.ASDR, 
                   Sweden = S.ASDR, World = W.ASDR) 

## Printing results
ASDR # Hard to analize!

## Plotting results
barplot(as.matrix(t(ASDR[,1:3])), beside = TRUE, 
        col = c("firebrick3", "deepskyblue3", "darkolivegreen3"),
        legend.text = colnames(ASDR),  # Legend matches column names
        args.legend = list(x = "topleft"),
        main = "Comparison of Age-Specific Death Rate",
        ylab = "ASDR", xlab = "Age Group",
        names.arg = rownames(ASDR))  # Labels for each category

## OBSERVATIONS:
# 1. The 2000-2005-period ASDRs of Kenya are significantly higher those of Sweden
# 2. The ASDR-gap Kenya vs Sweden is specially high for younger age groups

####
## Q5: Counterfactual CDR
####

S.p <- sweden$py[i2] / sum(sweden$py[i2]) # Sweden Age-Specific Person Year Proportion
cfK.CDR <- sum(K.ASDR * S.p)
## Compare factual and counterfactual CDRs.
K.CDR[[2]]; cfK.CDR

## OBSERVATIONS:
# 1. Kenya’s counterfactual CDR (using Sweden’s age distribution) is 0.023, 
#    compared to its factual CDR of 0.01 (more than twice). 
# 2. Therefore, age distribution affects the comparison of CDRs between countries.
\end{verbatim}

\subsection*{Summary of Results}
\begin{itemize}
    \item \textbf{Crude Birth Rates (CBR):} 
    \begin{itemize}
        \item Kenya’s CBR declined from 0.052 to 0.039 between 1950–1955 and 2005–2010.
        \item Sweden’s CBR also declined, but remained far more stable, going from 0.015 to 0.012 in between periods.
        \item The world’s CBR also declined, from 0.037 to 0.020.
    \end{itemize}

    \item \textbf{Total Fertility Rates (TFR):}
    \begin{itemize}
        \item Kenya’s TFR declined significantly, from 7.6 to 4.9.
        \item Sweden’s TFR decreased from 2.2 to 1.9.
        \item The world’s TFR dropped by half, from 5.0 to 2.5.
    \end{itemize}

    \item \textbf{Crude Death Rates (CDR):}
    \begin{itemize}
        \item Kenya’s CDR decreased from 0.024 to 0.016.
        \item Sweden’s CDR remained relatively constant at 0.01.
        \item The world’s CDR fell from 0.019 to 0.008.
        \item \textbf{Puzzling finding:} The CDR of Kenya in 2000-2005 is about the same as Sweden! This might be because the CDR does not take into account the age composition of a population. The next question explores this issue.
    \end{itemize}

    \item \textbf{Age-Specific Death Rates (ASDR):}
    \begin{itemize}
        \item The ASDR of Kenya is significantly higher than the ones of Sweden for each age group in the period 2000-2005
        \item In the same period, the ASDR-gap between Kenya and Sweden is specially higher for younger age groups
    \end{itemize}

    \item \textbf{Counterfactual CDR for Kenya:}
    \begin{itemize}
        \item Kenya’s counterfactual CDR, using Sweden’s age distribution, is 0.023, compared to its factual CDR of 0.01 (More than twice). 
        \item This demonstrates that age distribution affects the comparison of CDRs between countries.
    \end{itemize}
\end{itemize}
}

\end{document}